%% cover-letter.tex
%% Cover letter for submission to ACM Computing Surveys (CSUR)
%%
\documentclass[11pt]{letter}
\usepackage[margin=1in]{geometry}
\usepackage{hyperref}

\signature{Hameem Mahdi\\
Research Scientist\\
Mayo Clinic\\
Rochester, MN 55905, USA\\
\href{mailto:mahdi.hameem@mayo.edu}{mahdi.hameem@mayo.edu}\\
ORCID: \href{https://orcid.org/0009-0007-0005-3080}{0009-0007-0005-3080}}

\address{Hameem Mahdi\\Mayo Clinic\\Rochester, MN 55905, USA}

\begin{document}

\begin{letter}{%
  Editor-in-Chief\\
  ACM Computing Surveys (CSUR)\\
  Association for Computing Machinery}

\opening{Dear Editor-in-Chief,}

I am writing to submit the manuscript entitled \textbf{``A Unified Taxonomy of 19 AI System Types: Architecture, Capability, and Deployment Analysis for the Modern AI Landscape''} for consideration as a survey article in \textit{ACM Computing Surveys (CSUR)}.

\textbf{Summary.} This paper presents a comprehensive taxonomy of 19 distinct AI system types, organised into six functional clusters and analysed through three systematic dimensions: architectural design, capability spectrum, and deployment characteristics. The framework provides a structured vocabulary and analytical lens for understanding the increasingly diverse AI landscape, addressing a gap in the literature where existing classifications tend to focus narrowly on learning paradigms or application domains.

\textbf{Fit for ACM Computing Surveys.} This work is particularly well-suited for CSUR for the following reasons:

\begin{enumerate}
  \item \textbf{Survey scope alignment.} CSUR publishes comprehensive surveys and tutorials that cover the full spectrum of computing. Our taxonomy spans 19 system types across all major AI paradigms---from symbolic reasoning and probabilistic inference to generative models, agentic systems, and embodied AI---providing exactly the kind of broad, integrative coverage that defines CSUR's mission.

  \item \textbf{Unifying framework.} Unlike surveys that focus on individual subfields (e.g., XAI taxonomies, RL surveys, generative model typologies), this work provides a cross-cutting framework that reveals interconnections, convergence trends, and capability gaps across the full AI landscape. The three-dimensional analysis model (architecture, capability, deployment) offers a reusable analytical structure absent from existing classifications.

  \item \textbf{Rigorous validation.} The taxonomy is validated through four complementary empirical methods: (i) bibliometric analysis of 15,427 publications across Scopus, Web of Science, IEEE Xplore, and ACM Digital Library; (ii) an expert survey of 47 AI practitioners across 14 countries; (iii) mapping of 25 production AI systems; and (iv) industry adoption data from Gartner, McKinsey, and IDC reports.

  \item \textbf{Timeliness.} The rapid proliferation of AI system types---particularly agentic AI, multimodal models, and privacy-preserving systems---makes a unifying framework urgently needed. Our analysis shows Agentic AI publications growing at 89.1\% CAGR, yet no comprehensive taxonomy exists that situates these alongside established paradigms.

  \item \textbf{Practical impact.} Beyond classification, the paper provides actionable guidance for practitioners through a decision workflow for system selection, deployment trade-off analysis, and integration strategies. CSUR's broad readership of researchers, practitioners, and educators will benefit from this practical orientation.
\end{enumerate}

\textbf{Supplementary materials.} The submission includes supplementary materials containing the full expert survey instrument, detailed literature search methodology, complete case-study mapping (25 systems), and detailed survey results with demographic breakdowns.

\textbf{Preprint.} A preprint of this manuscript is publicly available on aiXiv (Paper ID: \texttt{aixiv.260413.000007}).

\textbf{Declarations.} This manuscript has not been published previously and is not under consideration at any other journal. There are no conflicts of interest to declare. The expert survey was conducted with appropriate ethical oversight and all responses were anonymised.

I believe this contribution will be of significant interest to CSUR's readership and look forward to your consideration.

\closing{Sincerely,}

\end{letter}
\end{document}
